
\section{\BookRecordIndex{T0356}辛涡旋方程}\label{sec:12-7}

本节讨论另一项与辛商的\BookRecordIndex{T0229}量子上同调密切相关的新进展。其出发点是把 \BookRecordIndex{T0003}$J$-全纯曲线的概念
推广到带哈密顿 $G$-作用的\BookRecordIndex{T0358}辛流形 $M$。新不变量由黎曼曲面 $\Sigma$ 上一种涡旋方程的解
构造；未知量的偶对由主 $G$-丛 $P\to\Sigma$ 上的一个联络，以及以 $M$ 为纤维的相伴丛
的一个截面组成。\BookRecordIndex{T0356}辛涡旋方程由 Salamon \srcite{64} 和 Ignasi Mundet i Riera
\srcite{294,295} 各自独立发现。当所论\BookRecordIndex{T0358}辛流形为 $\mathbb C^n$、$G$ 是
$\operatorname{U}(n)$ 的子群时，物理文献早已把这些方程称为规范化 sigma 模型。下面介绍
基本设置，并说明迄今在这个问题上得到的一些结果。

设 $(M,\omega)$ 是带紧 Lie 群 $G$ 的哈密顿作用的\BookRecordIndex{T0358}辛流形。这个作用由等变\BookRecordIndex{T0204}矩映射
\[
  \mu:M\longrightarrow\mathfrak g^*
\]
生成，其中 $\mathfrak g:=\operatorname{Lie}(G)$，而 $\mathfrak g^*$ 是其对偶。这意味着
元素 $\xi\in\mathfrak g$ 的无穷小作用 $X_\xi\in\operatorname{Vect}(M)$ 由
\[
  \iota(X_\xi)\omega=dH_\xi,
  \qquad H_\xi:=\langle\mu,\xi\rangle
\]
给出。令 $(\Sigma,j_\Sigma,\mathrm{dvol}_\Sigma)$ 为闭黎曼曲面，$P\to\Sigma$ 为
主 $G$-丛。\BookRecordIndex{T0356}辛涡旋方程的未知量是偶对 $(u,A)$，其中 $u:P\to M$ 是等变映射，
$A\in\mathcal A(P)$ 是联络；方程依赖参数 $\tau$，形式为
\begin{equation}\label{eq:12-7-1}
  \overline\partial_{J,A}(u)=0,
  \qquad *F_A+\mu(u)=\tau.
\end{equation}
这里 $J=\{J_z\}_{z\in\Sigma}$ 是 $M$ 上一族光滑的、$G$-不变且与 $\omega$ \BookRecordIndex{T0348}相容的\BookRecordIndex{T0088}殆复
结构。把联络看成等变的李代数值一形式 $A\in\Omega^1(P,\mathfrak g)$，也就是说
\[
  A_{pg}(vg)=g^{-1}A_p(v)g,
  \qquad v\in T_pP,\quad g\in G,
\]
并且它把竖直切空间与 $\mathfrak g$ 识别，即
\[
  A_p(p\xi)=\xi,
  \qquad p\in P,\quad \xi\in\mathfrak g.
\]
$u$ 关于 $A$ 的协变导数，是如下 $1$-形式：
\[
  d_Au\in\Omega^1(P,u^*TM),
  \qquad d_Au:=du+X_A(u).
\]
这个一形式等变且水平，因而下降为 $\Sigma$ 上取值于商丛 $u^*TM/G\to\Sigma$ 的一形式。
$M$ 上的\BookRecordIndex{T0088}殆复结构在这个向量丛上诱导复结构，而 $\overline\partial_{J,A}$ 表示 $d_Au$ 关于
这个复结构的复反线性部分。式~\eqref{eq:12-7-1} 中第二个方程把曲率
\[
  F_A:=dA+\frac12[A\wedge A]
  \in\Omega^2_{\mathrm{ad}}(P,\mathfrak g)
  =\Omega^2(\Sigma,\mathfrak g_P)
\]
与\BookRecordIndex{T0204}矩映射耦合。这里
\[
  *:\Omega^2(\Sigma,\mathfrak g_P)
  \longrightarrow\Omega^0(\Sigma,\mathfrak g_P^*)
\]
是 \BookRecordIndex{T0026}Hodge 星算子与同构 $\mathfrak g\to\mathfrak g^*$ 的复合，后一个同构由
$\mathfrak g$ 上的不变内积诱导。假设 $\tau\in Z(\mathfrak g^*)$ 属于中心，并且是\BookRecordIndex{T0204}矩映射
$\mu$ 的\BookRecordIndex{T0409}正则值。需要强调，$\Sigma$ 上的体积形式进入了~\eqref{eq:12-7-1} 的第二个方程；
因此，当 $J_z=J$ 与 $z$ 无关时，解空间在 $\Sigma$ 的复自同构群下不具有不变性。

规范群 $\mathcal G=\mathcal G(P)$ 由所有关于 $G$ 在自身上的共轭作用等变的光滑映射
$g:P\to G$ 组成。它在 $C_G^\infty(P,M)\times\mathcal A(P)$ 上有自然作用
\[
  g^*(u,A):=\bigl(g^{-1}u,g^{-1}dg+g^{-1}Ag\bigr),
\]
而这个作用保持式~\eqref{eq:12-7-1} 的解空间。

式~\eqref{eq:12-7-1} 的解是能量泛函
\[
  E(u,A):=\frac12\int_\Sigma
  \left(|d_Au|^2+|F_A|^2+|\mu(u)-\tau|^2\right)
  \mathrm{dvol}_\Sigma
\]
的绝对极小值。对式~\eqref{eq:12-7-1} 的解，能量由拓扑不变量给出：
\[
\begin{aligned}
  E(u,A)
  &=\int_\Sigma\left((d_Au)^*\omega
  -\langle\mu(u)-\tau,F_A\rangle\right)\\
  &=\int_\Sigma\left(u^*\omega
  -d\langle\mu(u)-\tau,A\rangle\right).
\end{aligned}
\]
这里两个被积式相等，并且都下降为 $\Sigma$ 上的二形式。这个积分可以解释为等变上同调类
\[
  [\omega-\mu+\tau]\in H_G^2(M)
\]
在等变同调类 $[u]\in H_2^G(M)$ 上的取值。（若 $\theta:P\to EG$ 是等变分类映射，则
$[u]$ 是基本类 $[\Sigma]\in H_2(\Sigma)$ 在如下映射下的像：这个映射
\[
  \Sigma\longrightarrow M_G:=M\times_GEG
\]
由 $u\times\theta:P\to M\times EG$ 诱导。）

\begin{remark}\label{rem:12-7-1}
\begin{enumerate}[label=(\roman*)]
\item 在 $\Sigma$ 上的局部全纯坐标 $s+it\in U$ 及 $P$ 的一个局部平凡化中，联络形如
\[
  A=\Phi\,ds+\Psi\,dt,
\]
其中 $\Phi,\Psi:U\to\mathfrak g$，而 $u$ 是从 $U$ 到 $M$ 的光滑映射。此时涡旋方程形如
\begin{equation}\label{eq:12-7-2}
\begin{aligned}
  \partial_su+X_\Phi(u)
  +J_{s,t}(u)\bigl(\partial_tu+X_\Psi(u)\bigr)&=0,\\
  \lambda^{-2}*_\mathfrak g
  \bigl(\partial_s\Psi-\partial_t\Phi+[\Phi,\Psi]\bigr)
  +\mu(u)&=\tau,
\end{aligned}
\end{equation}
其中 $*_\mathfrak g:\mathfrak g\to\mathfrak g^*$ 表示由 $\mathfrak g$ 上不变内积诱导的
同构，而局部坐标中的体积形式为
$\mathrm{dvol}_\Sigma=\lambda^2\,ds\wedge dt$。

\item 当 $G=\{1\}$ 时，式~\eqref{eq:12-7-1} 的解就是 $M$ 中的 \BookRecordIndex{T0003}$J$-全纯曲线。

\item 通常的涡旋方程是特例 $M=\mathbb C$、$G=S^1$。这时，式~\eqref{eq:12-7-1}
的第一个方程断言 $u$ 是复线丛 $L\to\Sigma$ 的全纯截面；设该线丛次数为 $d$，则在这种
情形，式~\eqref{eq:12-7-1} 的解模\BookRecordIndex{T0153}规范等价所得的\BookRecordIndex{T0246}模空间，自然等同于 $\Sigma$ 的 $d$ 重
对称积（见 Garcia-Prada \srcite{136}）。

\item 当 $M=\mathbb C^2$、$G=\operatorname{U}(2)$ 时，式~\eqref{eq:12-7-1} 的解\BookRecordIndex{T0246}模空间
是 \BookRecordIndex{T0011}Bradlow 偶对空间，即由 $\Sigma$ 上秩二全纯丛和一个全纯截面组成的稳定偶对之空间
\srcite{47}。此时参数 $\tau$ 决定稳定性条件。

\item 若 $M$ 是黎曼曲面 $S$ 上 $\operatorname{SU}(2)$-联络的无限维流形，而 $G$ 是
$S$ 上 $\operatorname{SU}(2)$-规范变换所成的空间，则式~\eqref{eq:12-7-1} 的解是乘积
$\Sigma\times S$ 上的反自对偶瞬子。

\item 若取 $M$ 为满足 $\overline\partial_A\Theta=0$ 的偶对
\[
  (\Theta,A)\in\Omega^0(S,E)\times\mathcal A(E)
\]
所成的空间，并取 $G=C^\infty(S,S^1)$，那么 $\Sigma\times S$ 上的
Seiberg--Witten 方程也作为式~\eqref{eq:12-7-1} 的特例出现；这里 $E\to S$ 是复线丛。
\end{enumerate}
\end{remark}

为了从式~\eqref{eq:12-7-1} 的解得到不变量，必须施加以下条件。当 $G=\{1\}$ 时，这些
条件等价于定义~\ref{def:9-2-6} 中的\BookRecordIndex{T0344}无穷远处凸性。

\begin{description}
\item[\normalfont\textbf{(H1)}]
\BookRecordIndex{T0204}矩映射
$\mu:M\to\mathfrak g^*$
\sourcefootnote{原书在这里把\BookRecordIndex{T0204}矩映射的陪域印作 $\mathfrak g$；但本节开头的定义以及同一
条件中的 $\tau\in Z(\mathfrak g^*)$ 均要求陪域为 $\mathfrak g^*$，故补正。}
是\BookRecordIndex{T0277}恰当映射，并且 $(M,\omega,\mu)$ 在如下意义下于\BookRecordIndex{T0344}无穷远处凸：存在
$G$-不变、与 $\omega$ \BookRecordIndex{T0348}相容的\BookRecordIndex{T0088}殆复结构 $J\in\mathcal J_G(M,\omega)$，一个 $G$-不变的
恰当光滑函数
\[
  f:M\longrightarrow[0,\infty),
\]
以及对每个 $\tau\in Z(\mathfrak g^*)$ 的紧集 $K_\tau\subset M$，使得
\[
  \omega_f(v,J(x)v)\geq0,
  \qquad
  df(x)J(x)X_{\mu(x)-\tau}\geq0
\]
对每个 $x\in M\setminus K_\tau$ 和每个 $v\in T_xM$ 成立。这里
\[
  \omega_f:=-d(df\circ J).
\]
\end{description}

要让式~\eqref{eq:12-7-1} 的解获得任何\BookRecordIndex{T0192}紧性，都需要假设 (H1)；没有它，便不能期待任何
有意义的结果。最简单却又非常有用的一类例子还满足

\begin{description}
\item[\normalfont\textbf{(H2)}]
$(M,\omega)$ \BookRecordIndex{T0355}辛非球面。
\end{description}

条件 (H1) 与 (H2) 合在一起迫使环境流形 $M$ 非紧。不过，它们蕴含：对每个
$B\in H_2^G(M)$，式~\eqref{eq:12-7-1} 中表示等变同调类 $[u]=B$ 的解，模基\BookRecordIndex{T0153}规范等价
所得的\BookRecordIndex{T0246}模空间 $\mathcal M_{B,\Sigma}(\tau)$ 是紧的（见 \srcite{64,67}）。由恰当\BookRecordIndex{T0204}矩映射
生成的 $\mathbb C^n$ 上每个线性作用都满足这两个假设。此外，只要
$\tau\in Z(\mathfrak g^*)$ 是\BookRecordIndex{T0204}矩映射的中心\BookRecordIndex{T0409}正则值且 $\Sigma$ 的体积充分大，这个\BookRecordIndex{T0246}模空间
就是正则的；这里正则是指规范群的作用具有有限稳定子。存在一个 $G$-等变\BookRecordIndex{T0283}求值映射
\[
  \operatorname{ev}:\mathcal M_{B,\Sigma}(\tau)\longrightarrow M,
\]
而参数 $\tau$ 处的不变量形如群同态
\[
  \operatorname{GGW}^{M,\mu-\tau}_{B,\Sigma}:
  H_G^*(M;\mathbb Q)\longrightarrow\mathbb Q.
\]
粗略地说，这个同态通过在商 $\mathcal M_{B,\Sigma}(\tau)/G$ 上“积分”等变上同调类
$a\in H_G^*(M;\mathbb Q)$ 的拉回来定义：
\begin{equation}\label{eq:12-7-3}
  \operatorname{GGW}^{M,\mu-\tau}_{B,\Sigma}(a)
  :=\int_{\mathcal M_{B,\Sigma}(\tau)/G}\operatorname{ev}^*a.
\end{equation}
这些就是\emph{\BookRecordIndex{T0154}规范化格罗莫夫--威滕不变量}。积分的严格定义不需要\BookRecordIndex{T0166}横截性；可以用一个
相伴 $G$-模问题的等变 Euler 类来理解它（见 Cieliebak--Mundet--Salamon
\srcite{66}），细节由 Cieliebak--Gaio--Mundet--Salamon 在 \srcite{67} 中给出。这个构造
还有若干细化版本，其中会改变\BookRecordIndex{T0066}标记点，或从联络空间拉回等变上同调类；这里不作讨论。

与第~\ref{sec:9-1} 节一样，可以用这些不变量证明时间依赖、$G$-等变哈密顿微分方程的
相对周期轨道存在，也就是证明存在模 $G$-作用闭合的解（见
Cieliebak--Gaio--Mundet--Salamon \srcite{67}）。另一个应用是：对适当流形 $M$，
Seiberg--Witten 不变量与\BookRecordIndex{T0154}规范化格罗莫夫--威滕不变量相对应（见
Cieliebak--Gaio--Mundet--Salamon \srcite{67} 和 Okonek--Teleman \srcite{309}）。这
重新得到了直纹曲面的 Seiberg--Witten 不变量之计算。事实上，假设 (H1--2) 下\BookRecordIndex{T0246}模空间紧，
所以涡旋不变量适于计算。一个有趣特征是它依赖参数 $\tau$。Cieliebak 与 Salamon 在
\srcite{69} 中对 $G$ 为环面的情形得到一个\BookRecordIndex{T0078}穿墙公式，描述当 $\tau$ 从一个\BookRecordIndex{T0280}腔跨入另一个\BookRecordIndex{T0280}腔
时不变量如何变化；他们还借此证明了单调\BookRecordIndex{T0169}环面流形的 Batyrev 公式。其证明的关键成分，是
\BookRecordIndex{T0154}规范化格罗莫夫--威滕不变量与辛商的\BookRecordIndex{T0143}格罗莫夫--威滕不变量之间的对应；下面解释这一点。

为简单起见，假设 $\tau=0$。要得到光滑商，还必须施加第三个条件：

\begin{description}
\item[\normalfont\textbf{(H3)}]
群 $G$ 在 $\mu^{-1}(0)$ 上自由作用。
\end{description}

在这个假设下，商
\[
  \overline M:=M/\!/G:=\mu^{-1}(0)/G
\]
是紧\BookRecordIndex{T0358}辛流形。\BookRecordIndex{T0003}$J$-全纯曲线 $\overline u:\Sigma\to\overline M$ 可以由满足
\begin{equation}\label{eq:12-7-4}
  \overline\partial_{J,A}(u)=0,
  \qquad \mu(u)=0
\end{equation}
的偶对
\[
  (u,A)\in C_G^\infty(P,M)\times\mathcal A(P)
\]
表示。要在式~\eqref{eq:12-7-1} 与~\eqref{eq:12-7-4} 的解之间建立对应，基本思路是研究
方程
\begin{equation}\label{eq:12-7-5}
  \overline\partial_{J,A}(u)=0,
  \qquad *F_A+\varepsilon^{-2}\mu(u)=0
\end{equation}
在 $\varepsilon$ 很小时的极限。Gaio 与 Salamon 在商单调的情形于 \srcite{134,135}
中完成了这一纲领，得到如下结果。

\begin{theorem}[Gaio--Salamon]\label{thm:12-7-2}
假设 (H1--3) 成立，并设 $c_1(T\overline M)$ 是 $[\overline\omega]$ 的正倍数。令
$\overline B\in H_2(\overline M;\mathbb Z)$，并令
$a_i\in H_G^*(M)$ 满足
\[
  \deg(a_i)<2N,
\]
其中 $N$ 是 $\overline M$ 的最小 Chern 数。以 $B\in H_2^G(M)$ 表示
$\overline B$ 在同态 $H_2(\overline M)\to H_2^G(M)$ 下的像，以
\[
  \overline a_i\in H^*(\overline M)
\]
表示 $a_i$ 在 Kirwan 同态
\[
  \kappa:H_G^*(M)\longrightarrow H^*(\overline M)
\]
下的像。那么
\[
  \operatorname{GW}^{\overline M}_{\overline B,\Sigma}
  (\overline a_1,\ldots,\overline a_k)
  =\operatorname{GGW}^{M,\mu}_{B,\Sigma}
  (a_1\smile\cdots\smile a_k).
\]
\end{theorem}

这里的\BookRecordIndex{T0143}格罗莫夫--威滕不变量使用\BookRecordIndex{T0157}哈密顿扰动、\BookRecordIndex{T0150}固定标记点和 $\Sigma$ 上的固定复结构来理解
（如第~\ref{ch:hamiltonian-perturbations} 章）。定理~\ref{thm:12-7-2} 的证明依赖于方程
\eqref{eq:12-7-5} 在 $\varepsilon\to0$ 时的绝热极限分析。这涉及颇为精细的估计，而证明
让人联想到 Dostoglou--Salamon \srcite{92} 对 Atiyah--Floer 猜想的证明。结合
定理~\ref{thm:12-7-2} 与习题~\ref{ex:11-1-18} 中把\BookRecordIndex{T0143}格罗莫夫--威滕不变量解释为从
$\overline M$ 的\BookRecordIndex{T0229}量子上同调到 $\mathbb Z$ 的同态，就得到下面的推论。

\begin{corollary}[Gaio--Salamon]\label{cor:12-7-3}
设 $(M,\omega,\mu)$ 如定理~\ref{thm:12-7-2}，并设 $\Lambda$ 是
定义~\ref{def:11-1-3} 意义下 $\overline M$ 的一个\BookRecordIndex{T0231}量子系数环。再假设 $M$ 的等变上同调
由次数小于 $2N$ 的类生成。那么存在唯一的环同态
\begin{equation}\label{eq:12-7-6}
  Q\kappa:H_G^*(M;\Lambda)
  \longrightarrow\operatorname{QH}^*(\overline M;\Lambda),
\end{equation}
它在次数小于 $2N$ 的类上与 Kirwan 同态
\[
  \kappa:H_G^*(M;\Lambda)\longrightarrow H^*(\overline M;\Lambda)
\]
一致。此外，$Q\kappa$ 满射，而且
\[
  \operatorname{GW}^{\overline M}_{\overline B,\Sigma}(Q\kappa(a))
  =\operatorname{GGW}^{M,\mu}_{B,\Sigma}(a)
\]
对每个 $\Sigma$、每个 $\overline B\in H_2(\overline M;\mathbb Z)$ 和每个
$a\in H_G^*(M)$ 成立。
\end{corollary}

由注~\ref{rem:11-3-13}，第~\ref{subsec:11-3-1}、\ref{subsec:11-3-2}、
\ref{subsec:11-3-3} 小节讨论的单调\BookRecordIndex{T0169}环面流形、\BookRecordIndex{T0274}旗流形和格拉斯曼流形的\BookRecordIndex{T0229}量子上同调环，都给出
推论~\ref{cor:12-7-3} 中环同态的例子，其中
\[
  M=\mathbb C^n,
  \qquad H_G^*(M;\Lambda)\cong H^*(BG;\Lambda).
\]

\begin{remark}[量子 Kirwan 同态]\label{rem:12-7-4}
\begin{enumerate}[label=(\roman*)]
\item 若削弱假设 (H3)，只假设零是\BookRecordIndex{T0204}矩映射的\BookRecordIndex{T0409}正则值，那么商 $\overline M$ 是辛轨形。
轨形的\BookRecordIndex{T0143}格罗莫夫--威滕不变量是一项新课题（见 Chen--Ruan \srcite{62}）。可以预期，
定理~\ref{thm:12-7-2} 以及推论~\ref{cor:12-7-3} 会相当直接地推广到商为单调辛轨形的
情形。

\item 正如 Gaio--Salamon \srcite{135} 所猜想的，推论~\ref{cor:12-7-3} 还应推广到非单调
情形；不过证明需要细化 \srcite{135} 中的分析，并涉及\BookRecordIndex{T0356}辛涡旋方程解的\BookRecordIndex{T0256}粘合定理。其中一个
成分，是通过复平面 $\Sigma=\mathbb C$ 上\BookRecordIndex{T0356}辛涡旋方程~\eqref{eq:12-7-2} 的有限能量解
（取 $\lambda\equiv1$），直接给出同态 $Q\kappa$ 的几何定义；比较 Ziltener
\srcite{430}。

\item 去掉假设 (H2) 则是另一回事；这个假设断言环境流形\BookRecordIndex{T0355}辛非球面。这样做需要在 $M$ 的
等变上同调中引入量子形变项。近年已有若干研究者在这些猜想上取得大量进展；参见
Ziltener \srcite{431,432}、Gonzalez--Woodward \srcite{156,157}、Ott \srcite{319}、
Nguyen--Woodward--Ziltener \srcite{299}。在这些文献中，式~\eqref{eq:12-7-6} 的同态
$Q\kappa$ 及其猜想推广称为\emph{量子 Kirwan 同态}。
\end{enumerate}
\end{remark}

\begin{remark}[\BookRecordIndex{T0203}矩弗洛尔同调]\label{rem:12-7-5}
\BookRecordIndex{T0132}弗洛尔同调有一个与\BookRecordIndex{T0356}辛涡旋方程相关的版本，正如标准辛\BookRecordIndex{T0132}弗洛尔同调与 \BookRecordIndex{T0003}$J$-全纯曲线相关。
在这个理论中，作用泛函定义在 $M$ 与 Lie 代数的乘积中环路
\[
  (x,\eta):\mathbb R/\mathbb Z\longrightarrow M\times\mathfrak g
\]
所成的空间上，形式为
\[
  \mathcal A_{H,\mu}(x,\eta)
  =-\int_Bu^*\omega
  +\int_0^1\left(
  \langle\mu(x(t)),\eta(t)\rangle-H_t(x(t))
  \right)dt.
\]
这里 $H_t=H_{t+1}:M\to\mathbb R$ 是一族光滑的、$1$-周期且 $G$-不变的哈密顿函数；并且
和第~\ref{sec:12-1} 节一样，$u:B\to M$ 是满足
$u(e^{2\pi it})=x(t)$ 的光滑映射。这个泛函的\BookRecordIndex{T0233}临界点模\BookRecordIndex{T0153}规范等价所得的类，可以等同于
$\mu^{-1}(0)$ 中哈密顿系统相对周期轨道的等价类。

Frauenfelder 在 \srcite{121,122} 中使用这个弗洛尔理论的拉格朗日版本，证明了带
$G$-作用的\BookRecordIndex{T0355}辛非球面流形之辛商的 Arnold--Givental 猜想。这个猜想断言：当所论\BookRecordIndex{T0219}拉格朗日
子流形是某个反辛对合的不动点集时，它与自身任一个哈密顿形变的交点个数有一个下界。
Fukaya--Oh--Ohta--Ono \srcite{128} 较早给出了单调情形的证明，不过不带商空间假设。
\end{remark}
